%==============================================================================
%  CHAPTER 2 — PLANNING AND CAPTURE OF NEEDS
%==============================================================================
\chapter{Planning and Capture of Needs}
\label{chap:requirements}

\section*{Introduction}
This chapter captures what the platform must do and how the work was organized. It
analyzes the needs (actors, functional and non-functional requirements, and the
use-case model), presents the project management with Scrum (the team, the product
backlog, and the sprint plan), and finally describes the application architecture
and the overall architecture of the project.

% ---------------------------------------------------------------------------
\section{Needs Analysis}
\label{sec:requirements}

\subsection{Identification of Actors}
An actor is any entity that interacts with the system from the outside. Four
actors were identified.

\paragraph{Platform user (data scientist).}
The primary actor and the reason the platform exists. After authenticating, the
user organizes work into projects, uploads code and data, triggers training runs
and follows their live logs, inspects and compares experiments, registers and
promotes model versions, deploys them as live services, tests predictions from
the interface, and manages the \gls{api} keys that expose a deployment publicly.
The platform assumes no infrastructure knowledge from this actor: every cluster
interaction happens behind a button.

\paragraph{External API client.}
A third-party application, a website, a mobile application, or any backend
script, that calls a deployed model's public prediction endpoint using an
\gls{api} key. This actor has no platform account and never sees the interface;
its entire contract with the system is one authenticated \gls{http} call.

\paragraph{Administrator.}
A privileged user who oversees the whole platform: manages user accounts
(activate, deactivate, promote, demote, purge, delete), views platform-wide
statistics and cluster health, and inspects every user's projects, pipelines, and
deployments. An administrator inherits all capabilities of a platform user.

\paragraph{System (automated actors).}
The training pipeline and the model server act on the user's behalf once
triggered: the pipeline downloads code, runs it, and reports results to the
tracking server; the model server loads a model artifact and answers prediction
requests. They are secondary actors that appear in the sequence diagrams of the
sprint chapters.

\subsection{Functional Specification}
\begin{itemize}[leftmargin=1.4em]
  \item \textbf{Authentication.} Sign up, sign in, and session refresh using
        \gls{jwt} tokens.
  \item \textbf{Project management.} Create, list, update, and delete projects;
        each project is bound to its own experiment-tracking namespace.
  \item \textbf{Upload.} Upload code (\texttt{.py}, \texttt{.ipynb}, \texttt{.zip})
        and datasets (\texttt{.csv}) through the browser.
  \item \textbf{Training.} Trigger a training pipeline that executes the uploaded
        code on the cluster.
  \item \textbf{Experiment tracking.} Automatically capture metrics, parameters,
        and artifacts; browse and compare runs.
  \item \textbf{Model registry.} Register model versions and promote them between
        Staging, Production, and Archived stages.
  \item \textbf{Deployment.} Deploy a registered model as a live prediction
        service.
  \item \textbf{Prediction.} Send sample inputs and receive predictions, both from
        the \gls{ui} and from a public \gls{api} endpoint secured by per-deployment
        keys.
  \item \textbf{Monitoring.} View live cluster and per-pod resource metrics on a
        dashboard.
  \item \textbf{Administration.} Manage user accounts and roles, view
        platform-wide statistics and cluster health, and oversee all users'
        projects, pipelines, and deployments (administrator only).
\end{itemize}

\subsection{Non-Functional Requirements}
Beyond what the system does, the following qualities constrain how it does it.

\begin{itemize}[leftmargin=1.4em]
  \item \textbf{Security.} Authentication is stateless, based on signed \gls{jwt}
        tokens~\cite{rfc7519}, so any backend replica can validate a request.
        Passwords are stored only as bcrypt hashes, public \gls{api} keys only as
        \gls{sha}-256 hashes, in line with \gls{api} security good
        practice~\cite{owasp_api}. Every resource access is checked against the
        owner, and administrative endpoints require the administrator role.
  \item \textbf{Performance.} The backend is fully asynchronous, so slow external
        calls (cluster \gls{api}, object store, tracking server) never block other
        requests; the interface must remain responsive while a training run
        executes.
  \item \textbf{Scalability.} Each training job and each model server runs in its
        own resource-limited pod on Kubernetes, so workloads are isolated from
        one another and the cluster can grow by adding nodes.
  \item \textbf{Portability.} Everything is containerized and expressed as
        Kubernetes resources; the platform runs identically on a local KinD
        cluster and on \gls{aks}, which is how it was developed and deployed.
  \item \textbf{Reproducibility.} A run's code, parameters, metrics, and
        artifacts are captured automatically, so any result can be traced back to
        what produced it.
  \item \textbf{Usability.} The whole lifecycle is driven from a single, coherent
        web interface; a data scientist needs no kubectl, no YAML, and no cloud
        console.
  \item \textbf{Maintainability.} The backend is organized in layers (routers,
        services, models) with typed schemas at the boundaries, so each concern
        can evolve independently.
\end{itemize}

\subsection{Use-Case Model}
\label{sec:usecases}
Figure~\ref{fig:usecase-global} presents the global use-case diagram, summarizing
the interactions of the three human actors (platform user, administrator, and
external \gls{api} client) with the system.
\begin{figure}[H]\centering
  \includegraphics[width=0.9\textwidth,height=0.72\textheight,keepaspectratio]{usecase_global}
  \caption{Global use-case diagram of the platform}
  \label{fig:usecase-global}
\end{figure}

% ---------------------------------------------------------------------------
\section{Project Management with Scrum}
\label{sec:scrum-management}

\subsection{The Scrum Team}
The Scrum framework presented in Chapter~\ref{chap:context} was applied with the
following team:
\begin{itemize}[leftmargin=1.4em]
  \item \textbf{\gls{po}}: Mr.\ Amine Gonji (company supervisor, INSOMEA),
        responsible for prioritizing the product backlog and validating each
        increment.
  \item \textbf{\gls{sm}}: Mr.\ Ben Mardes Achref (academic supervisor),
        responsible for facilitating the process and removing impediments.
  \item \textbf{Development Team}: the student (author), responsible for analysis,
        design, implementation, and testing, with technical guidance from
        Mr.\ Benarous Abderrahmen.
\end{itemize}

\subsection{Product Backlog}
\label{sec:backlog}
The product backlog (Table~\ref{tab:backlog}) gathers all desired capabilities as
user stories, ordered by priority; story points estimate the relative effort.

\begin{table}[H]\centering
\caption{Product backlog}\label{tab:backlog}
\renewcommand{\arraystretch}{1.25}
\begin{tabularx}{\textwidth}{|c|X|c|c|}
\hline
\textbf{\#} & \textbf{User story} & \textbf{Prio.} & \textbf{Pts}\\
\hline
1  & As a user, I want to sign up and sign in so that my work is private.        & High & 5\\ \hline
2  & As a user, I want to create and manage projects to organize my work.        & High & 5\\ \hline
3  & As a user, I want to upload code and a dataset through the browser.         & High & 5\\ \hline
4  & As a user, I want to trigger a training run on the cluster.                 & High & 8\\ \hline
5  & As a user, I want metrics and artifacts captured automatically.             & High & 5\\ \hline
6  & As a user, I want to watch the live logs of a running pipeline.             & Med  & 3\\ \hline
7  & As a user, I want to browse and compare my experiment runs.                 & Med  & 3\\ \hline
8  & As a user, I want to register and version trained models.                   & High & 5\\ \hline
9  & As a user, I want to promote a model between Staging/Production/Archived.   & High & 3\\ \hline
10 & As a user, I want to deploy a registered model as a live service.           & High & 8\\ \hline
11 & As a user, I want to send inputs and get predictions from the UI.           & High & 5\\ \hline
12 & As a user, I want a dashboard of cluster and per-pod metrics.               & Med  & 5\\ \hline
13 & As a user, I want public API keys so external apps can call my model.       & Med  & 5\\ \hline
14 & As a user, I want clear diagnostics when an uploaded script fails.          & Med  & 5\\ \hline
15 & As an administrator, I want to manage users and oversee the platform.       & Med  & 5\\ \hline
\end{tabularx}
\end{table}

\subsection{Sprint Planning}
\label{sec:release-plan}
The backlog was delivered over four sprints, each producing a demonstrable
increment:

\begin{itemize}[leftmargin=1.4em]
  \item \textbf{Sprint~1, Foundation:} authentication and project management
        (stories 1--2).
  \item \textbf{Sprint~2, Training \& Tracking:} upload, training pipeline, live
        logs, and experiment tracking (stories 3--7).
  \item \textbf{Sprint~3, Registry \& Deployment:} model registry, promotion,
        deployment, and prediction (stories 8--11).
  \item \textbf{Sprint~4, Observability \& Public Serving:} metrics dashboard,
        public \gls{api} keys, diagnostics, and administration (stories 12--15).
\end{itemize}

This split is deliberate: each sprint delivers a complete vertical slice that the
previous one makes possible. Sprint~1 establishes identity and ownership, without
which nothing can be private; Sprint~2 produces tracked training runs, which are
the raw material of Sprint~3's registry and deployments; Sprint~4 then makes the
now-complete lifecycle observable, externally consumable, and administrable. The
story numbers in each sprint chapter's backlog trace back to
Table~\ref{tab:backlog}, so every implemented task is connected to a requirement.

The project ran over six months, from 12 January 2026, opening with a framing and
state-of-the-art phase, followed by the four sprints, with the report written in
parallel throughout. Table~\ref{tab:schedule} lists the phases with their dates and
durations.

\begin{table}[H]\centering
\renewcommand{\arraystretch}{1.35}
\caption{Project schedule}\label{tab:schedule}
\begin{tabularx}{\textwidth}{|>{\bfseries}l|c|c|X|}
\hline
\rowcolor{primary!15}
\textbf{Phase} & \textbf{Period} & \textbf{Duration} & \textbf{Main objective}\\
\hline
Framing \& SOTA & 12 Jan -- 01 Feb & 3 weeks & Framing, state of the art, and environment setup\\ \hline
Sprint 1 & 02 Feb -- 01 Mar & 4 weeks & Authentication and project management\\ \hline
Sprint 2 & 02 Mar -- 05 Apr & 5 weeks & Training pipeline and experiment tracking\\ \hline
Sprint 3 & 06 Apr -- 03 May & 4 weeks & Model registry, deployment, and prediction\\ \hline
Sprint 4 & 04 May -- 14 Jun & 6 weeks & Observability, public API, and administration\\ \hline
Report writing & 12 Jan -- 11 Jul & Continuous & Drafting and revising the report in parallel\\ \hline
\end{tabularx}
\end{table}

% ---------------------------------------------------------------------------
\section{Application Architecture}
\label{sec:app-architecture}
The platform adopts a \textbf{layered (n-tier) architecture} that enforces a clear
separation of concerns and keeps each responsibility independently testable and
replaceable:

\begin{itemize}[leftmargin=1.4em]
  \item \textbf{Presentation layer}, an Angular single-page application that runs
        in the browser and consumes the backend's \gls{rest} \gls{api}.
  \item \textbf{API layer}, FastAPI routers that validate requests, enforce
        authentication and ownership, and expose the platform's endpoints.
  \item \textbf{Service layer}, the business logic that orchestrates the external
        systems (pipeline submission, MLflow tracking, KServe deployment, object
        storage), isolating the routers from infrastructure details.
  \item \textbf{Data layer}, SQLAlchemy models mapped onto PostgreSQL, holding
        users, projects, pipeline runs, models, deployments, and API keys.
\end{itemize}

Figure~\ref{fig:layered-arch} shows these layers together with the cluster systems
that the service layer orchestrates.

\begin{figure}[H]\centering
\begin{tikzpicture}
  \node[layernode,minimum width=6.6cm] (pres) {Presentation --- Angular SPA};
  \node[layernode,minimum width=6.6cm,below=0.55cm of pres] (api) {API --- FastAPI routers};
  \node[layernode,minimum width=6.6cm,below=0.55cm of api] (svc) {Service --- business logic};
  \node[layernode,minimum width=6.6cm,below=0.55cm of svc] (data) {Data --- SQLAlchemy / PostgreSQL};
  \foreach \a/\b in {pres/api,api/svc,svc/data}{
    \draw[flowarrow] (\a) -- (\b);
  }
  \node[storenode,right=2.4cm of svc,yshift=1.7cm] (mlf) {
\includegraphics[height=0.4cm]{mlflow}\ MLflow};
  \node[storenode,right=2.4cm of svc,yshift=0.55cm] (mio) {
\includegraphics[height=0.38cm]{minio}\ MinIO};
  \node[storenode,right=2.4cm of svc,yshift=-0.6cm] (ksv) {
\includegraphics[height=0.38cm]{kserve}\ KServe};
  \node[storenode,right=2.4cm of svc,yshift=-1.75cm] (kfp) {
\includegraphics[height=0.38cm]{kubeflow}\ KFP / Argo};
  \foreach \t in {mlf,mio,ksv,kfp}{
    \draw[flowarrow,draw=diagamber!75] (svc.east) -- (\t.west);
  }
  \begin{scope}[on background layer]
    \node[clusterbox,fit=(mlf)(mio)(ksv)(kfp),
          label={[font=\scriptsize\bfseries,text=diagblue]above:{Cluster systems}}] {};
  \end{scope}
\end{tikzpicture}
\caption{Layered application architecture and the cluster systems it orchestrates}
\label{fig:layered-arch}
\end{figure}

This separation means, for example, that the deployment logic can talk to KServe
through a single service module without the \gls{api} layer knowing how an
\texttt{InferenceService} is created, and that the data layer can evolve
independently of the user interface.

The layered style was preferred over a microservices decomposition for a reason
of proportionality: the platform is operated by a small team and already
delegates its heavyweight concerns (pipeline execution, model serving, artifact
storage) to dedicated cluster components. Splitting the remaining business logic
into separate services would multiply deployment units and network boundaries
without adding isolation where it matters, which the per-pod execution model
already provides. The result is one well-structured backend image and one
frontend image, each independently deployable, in front of a set of specialized
open-source systems.

% ---------------------------------------------------------------------------
\section{Overall Project Architecture}
\label{sec:architecture}
At the infrastructure level, the platform is fully containerized and deployed on
Kubernetes. A browser-based \gls{spa} talks to the asynchronous backend, which
orchestrates the training pipeline, the experiment tracker, the object store, the
database, and the model servers; public traffic reaches the cluster through a
Cloudflare Tunnel. Figure~\ref{fig:architecture} shows the overall deployment.

\begin{figure}[H]\centering
  \includegraphics[width=0.9\textwidth,height=0.72\textheight,keepaspectratio]{architecture}
  \caption{Global deployment architecture of the platform}
  \label{fig:architecture}
\end{figure}

The request path is: \emph{browser $\rightarrow$ Cloudflare edge (\gls{tls})
$\rightarrow$ tunnel agent in the cluster $\rightarrow$ frontend service
$\rightarrow$ backend service $\rightarrow$ (database / object store / tracker /
pipeline / model servers)}.

\section*{Conclusion}
This chapter captured the platform's needs, organized the work with Scrum, the
team, a prioritized product backlog, and a four-sprint plan, and described the
application and overall project architecture. The following four chapters detail
the study and implementation of each sprint in turn.
